\subsection{The second existence theorem}

THEOREM 3. Let A be an integrally closed domain and A' a ring containing A and integral over A. Suppose that 0 is the only element of A which is a divisor of 0 in A'. Let p, q be two prime ideals of A such that q ⊃ p and q' a prime ideal of A' lying above q. Then there exists a prime ideal p' of A' lying above p and such that q' ⊃ p'.

Suppose first that A' is an integral domain. Let K, K' be the fields of fractions of A and A' respectively; let K'' be the algebraic closure of K' and A'' the integral closure of A in K''; then \(A \subset A' \subset A''\) . Let \(p''\) be a prime ideal of A'' lying above p (no. 1, Theorem 1), \(q''\) a prime ideal of A'' lying above q and such that \(p'' \subset q''\) (no. 1, Corollary 2 to Theorem 1) and finally \(q_{1}''\) a prime ideal of A'' lying above \(q'\) (no. 1, Theorem 1). By no. 3, Proposition 6 (i), there exists a K-automorphism \(\sigma\) of K'' such that \(\sigma \cdot q'' = q_{1}''\) . Then \(\sigma \cdot p''\) is a prime ideal of A'' lying above p such that \(\sigma \cdot p'' \subset q_{1}''\) and hence \(p' = A' \cap \sigma \cdot p''\) is a prime ideal of A' lying above p and contained in \(A' \cap q_{1}'' = q'\) .

We pass to the general case. As A is an integral domain and \(q'\) is prime, the subsets \(A - \{0\}\) and \(A' - q'\) of \(A'\) are multiplicative; then their product \(S = (A - \{0\})(A' - q')\) is a multiplicative subset of \(A'\) which does not contain 0 since the non-zero elements of A are not divisors of 0 in \(A'\) . Then there exists (Chapter II, §2, no. 5, Corollary 2 to Proposition 11) a prime ideal \(m'\) of \(A'\) disjoint from S, in other words such that \(m' \subset q'\) and \(m' \cap A = 0\) . Let h be the canonical homomorphism \(A' \to A'/m'\) . The restriction of h to A is injective and hence \(h(A)\) is integrally closed. As \(m' \subset q'\) , \(h(q')\) is a prime ideal of \(A'/m'\) lying above \(h(q)\) ; since \(A'/m'\) is an integral domain, the first part of the proof proves that there exists a prime ideal \(n'\) of \(A'/m'\) such that \(n' \cap h(A) = h(p)\) and \(h(q') \supset n'\) . The ideal \(p' = h^{-1}(n')\) is a prime ideal of \(A'\) and \(q' \supset p'\) , since \(q'\) contains the kernel of h. As \(n' \supset h(p)\) , \(p' \supset p\) . Finally, for \(x \in p' \cap A\) , \(h(x) \in n' \cap h(A) = h(p)\) and hence \(x \in p\) since the restriction of h to A is injective; hence \(p' \cap A = p\) .

COROLLARY. Under the hypotheses on A and A' of Theorem 3, let p be a prime ideal of A. The prime ideals of A' lying above p are the minimal elements of the set E of prime ideals of A' containing pA'.

A prime ideal of A' lying above p is minimal in E by virtue of no. 1, Corollary 1 to Proposition 1. Conversely, let q' be a minimal element of E. As q' ∩ A ⊃ p, Theorem 3 shows that there exists a prime ideal p' of A' lying above p such that q' ⊃ p'. As p' contains pA', the hypothesis made on q' implies that q' = p' and hence q' lies above p.

* Let V, V' be two affine algebraic varieties and f a morphism of V' to V such that \(f(V')\) is dense in V. Let A (resp. A') be the ring of functions regular on V (resp. V'); having been given f we can identify A with a subring of A'; suppose that A' is integral over A. Theorem 1 of no. 1 shows that for every irreducible subvariety W of V there exists an irreducible subvariety \(W'\) of \(V'\) such that \(f(W')\) is a dense subset of W; in particular, every point of V is the image of an irreducible subvariety of \(V'\) , which shows that f is surjective. Similarly, the restriction of f to every irreducible subvariety \(W'\) of \(V'\) maps \(W'\) onto an irreducible subvariety of V. Corollary 2 to Theorem 1, no. 1 shows that, if W and X are two irreducible subvarieties of V such that \(W \supset X\) and if \(W'\) is an irreducible subvariety of \(V'\) such that \(f(W') = W\) , then there exists an irreducible subvariety \(X'\) of \(V'\) contained in \(W'\) and such that \(f(X') = X\) .

If A is integrally closed, V is called normal. Theorem 3 shows that, if V is normal, if W and X are irreducible subvarieties of V such that \(W \supset X\) and if \(X'\) is an irreducible subvariety of \(V'\) such that \(f(\mathbf{X}') = \mathbf{X}\) , then there exists a subvariety \(W'\) of \(V'\) containing \(X'\) and such that \(f(W') = W\) . Finally, the Corollary to Theorem 3 shows that, if V is normal and W is an irreducible subvariety of V, the irreducible subvarieties \(W'\) of \(V'\) such that \(f(W') = W\) are just the irreducible components of \(f(W)\) .

